\clearpage
\section*{3. Lọc biểu cảm và state machine one-shot}
\addcontentsline{toc}{section}{3. Lọc biểu cảm và state machine one-shot}

\subsection*{3.1. Tại sao cần lọc theo thời gian?}

Output của mô hình cảm xúc có thể thay đổi giữa hai frame liên tiếp do ánh
sáng, góc mặt hoặc biểu cảm chuyển tiếp. Nếu mỗi nhãn được nối trực tiếp với
\path{/cmd_vel}, robot có thể đổi hướng liên tục. Module
\texttt{EmotionFilter} vì vậy tách ba khái niệm: observation mới từ transport,
ứng viên đang tích lũy thời gian và biểu cảm đã ổn định.

Một observation chỉ hợp lệ để làm mới heartbeat khi sequence tăng nghiêm ngặt
trong session hiện tại. Sequence trùng hoặc lùi sinh reason
\texttt{invalid\_sequence}; session mới sinh reason
\texttt{recognition\_restarted}, xóa candidate/stable cũ và buộc dừng. Đồng hồ
dùng \texttt{time.monotonic()} nên không phụ thuộc simulation time.

\subsection*{3.2. Chuẩn hóa và các ngưỡng}

Filter chuẩn hóa theo thứ tự an toàn sau:

\begin{enumerate}[leftmargin=1cm,itemsep=0.15em]
  \item \texttt{face\_count=0}: \texttt{unknown/no\_face};
  \item \texttt{face\_count} khác 1: \texttt{unknown/multiple\_faces};
  \item nhãn \texttt{disgust}: \texttt{unknown/unsupported\_emotion};
  \item nhãn ngoài contract: \texttt{unknown/unknown\_emotion};
  \item confidence dưới 0,55: \texttt{unknown/low\_confidence};
  \item còn lại: giữ nhãn điều khiển.
\end{enumerate}

Mất mặt, nhiều mặt, confidence thấp và nhãn không hỗ trợ tạo yêu cầu dừng ngay,
đồng thời vẫn cho \texttt{unknown} tích lũy đủ thời gian để re-arm. Các ngưỡng
mặc định được đặt trong \path{config/controller.yaml}:

\begin{center}
\begin{tabular}{@{}p{0.38\textwidth}p{0.18\textwidth}p{0.36\textwidth}@{}}
\toprule
\textbf{Tham số} & \textbf{Giá trị} & \textbf{Ý nghĩa} \\ \midrule
\texttt{confidence\_threshold} & 0,55 & Confidence tối thiểu cho một mặt. \\
\texttt{stable\_seconds} & 0,50~s & Thời gian ổn định của biểu cảm hành động. \\
\texttt{rearm\_seconds} & 0,75~s & Thời gian neutral/unknown để arm lại. \\
\texttt{observation\_timeout} & 0,75~s & Khoảng tối đa không có sequence mới. \\
\texttt{odom\_timeout} & 0,50~s & Tuổi tối đa của odometry. \\
\bottomrule
\end{tabular}
\end{center}

Khi nhãn chuẩn hóa đổi, candidate và thời điểm bắt đầu được đặt lại. Một nhãn
non-neutral trở thành stable sau 0,5~s liên tục. Neutral và unknown cần 0,75~s
vì chúng còn đóng vai trò xác nhận người dùng đã rời biểu cảm trước đó để hệ
thống được phép nhận một trigger mới.

\subsection*{3.3. State machine one-shot}

\begin{figure}[H]
\centering
\begin{tikzpicture}[
  state/.style={draw,rounded corners,minimum width=3.2cm,minimum height=1.1cm,
                align=center,fill=blue!6},
  arrow/.style={-{Latex[length=2.5mm]},thick},
  label/.style={align=center,font=\small}
]
  \node[state] (disarmed) {DISARMED\\chờ re-arm};
  \node[state,right=4.1cm of disarmed] (armed) {ARMED\\sẵn sàng nhận trigger};
  \node[state,below=2.4cm of armed] (executing) {EXECUTING\\đang chạy profile};

  \draw[arrow] (disarmed) -- node[label,above]
    {neutral/unknown\\ổn định 0,75 s} (armed);
  \draw[arrow] (armed) -- node[label,right]
    {biểu cảm hành động\\ổn định 0,5 s} (executing);
  \draw[arrow] (executing) -- node[label,below left]
    {hoàn tất, lỗi, timeout\\hoặc biểu cảm khác} (disarmed);
  \draw[arrow] (armed) to[bend left=40] node[label,above]
    {preview hoặc thiếu odom\\trigger vẫn bị tiêu thụ} (disarmed);
  \draw[arrow] (executing) to[bend left=28] node[label,right]
    {neutral/unknown\\hoặc stop reason} (armed);
\end{tikzpicture}
\caption{State machine một lần kích hoạt, không xếp hàng hành động.}
\end{figure}

Ánh xạ biểu cảm non-neutral được lưu cố định trong
\texttt{ACTION\_BY\_EMOTION}. Khi state là \texttt{ARMED} và filter báo một
stable emotion mới, state machine trả về đúng một \texttt{start\_action}. Khi
profile hoàn thành, state trở lại \texttt{DISARMED}; giữ nguyên nét mặt không
thể tự chạy lần nữa. Hệ thống không lưu action chờ, vì vậy một biểu cảm khác
trong lúc robot chạy chỉ hủy profile hiện tại sau khi nó ổn định, không nối
tiếp thành chuỗi ngoài ý muốn.

Neutral là lệnh dừng và re-arm. Unknown an toàn cũng có thể re-arm sau 0,75~s,
giúp người dùng chỉ cần rời khỏi camera giữa hai lần điều khiển. Payload lỗi,
session mới và timeout đều hủy chuyển động nhưng không tự kích hoạt hành động
khi dữ liệu quay trở lại.

\subsection*{3.4. Preview và odometry chưa sẵn sàng}

Trong preview, filter và state machine vẫn chạy để người dùng quan sát đúng
nhãn/stable state. Khi một biểu cảm đủ điều kiện kích hoạt, reason chuyển thành
\texttt{preview\_completed}; trigger bị tiêu thụ nhưng controller chỉ publish
zero Twist. Người dùng phải dừng phiên và truyền
\texttt{preview\_only:=false} để cho phép chuyển động.

Nếu motion được bật nhưng \path{/odom} chưa sẵn sàng, trigger bị từ chối với
reason \texttt{odometry\_unavailable} và cũng bị tiêu thụ. Khi odometry xuất
hiện sau đó, action cũ không được tự chạy; cần neutral/unknown re-arm rồi tạo
biểu cảm mới. Quy tắc này loại bỏ chuyển động bất ngờ sau khi simulator vừa
khởi động.

\reportfigure{05_preview_zero_cmd_vel.png}{Preview nhận trigger nhưng /cmd\_vel vẫn bằng zero; status hiển thị preview\_completed.}
